%%%%%%%%%%%%%%%%%%%%%%%%%%%%%%%%%%%%%%%%%%%%%%%%%%%%%%%%%%%%%%%%%%%%%%%%%%%%%%%
% Mutual information estimation
%%%%%%%%%%%%%%%%%%%%%%%%%%%%%%%%%%%%%%%%%%%%%%%%%%%%%%%%%%%%%%%%%%%%%%%%%%%%%%%

\subsection{Motivation and Problem Setting}
\label{sec:mi-est-motivation}

Mutual information (MI) is a fundamental measure of statistical dependence.
In representation learning it is often used as an objective:
maximizing $I(T;Y)$ encourages a representation $T$ to preserve information useful for predicting $Y$;
maximizing $I(T_1;T_2)$ under different ``views'' $T_1,T_2$ underlies contrastive learning;
and minimizing certain MI terms appears in the information bottleneck.

From Section~\ref{sec:entropy-mi},
\begin{equation}
    I(X;Y) = \KL{p_{X,Y}}{p_X p_Y} = \E{\log\frac{p(X,Y)}{p(X)p(Y)}}.
\end{equation}
The challenge is that $p(x,y)$ and the marginals are typically unknown.
Instead, we usually have only i.i.d. samples
\((x_i,y_i)\sim p_{X,Y}\).
Sampling from the product of marginals $p_Xp_Y$ is also easy: one can shuffle pairs (take $x_i$ with $y_j$ for $j\neq i$).

\paragraph{Density-ratio viewpoint.}
Define the density ratio
\begin{equation}
    r(x,y)\coloneqq \frac{p(x,y)}{p(x)p(y)}.
\end{equation}
Then $I(X;Y)=\Esub{p(x,y)}{\log r(X,Y)}$.
Most neural MI estimators learn either $\log r$ directly (a \emph{critic} function), or a classifier whose logit is related to $\log r$.

\subsection{Variational Lower Bounds on Mutual Information}
\label{sec:mi-est-bounds}

All bounds below are derived from variational characterizations of KL divergence.
We write expectations explicitly with respect to the joint $p_{X,Y}$ and the product of marginals $p_Xp_Y$.

\subsubsection*{1) Barber--Agakov (BA) bound}

Start from
\(I(X;Y)=H(X)-H(X\mid Y)\).
For any conditional distribution $q(x\mid y)$,
\begin{equation}
\label{eq:ba-bound}
\begin{aligned}
I(X;Y)
&= H(X) + \Esub{p(x,y)}{\log p(x\mid y)}\\
&= H(X) + \Esub{p(x,y)}{\log q(x\mid y)} + \Esub{p(y)}{\KL{p(x\mid y)}{q(x\mid y)}}\\
&\ge H(X) + \Esub{p(x,y)}{\log q(x\mid y)}.
\end{aligned}
\end{equation}
Thus
\begin{equation}
    I(X;Y) \ge I_{\mathrm{BA}}(q) \coloneqq H(X) + \Esub{p(x,y)}{\log q(x\mid y)}.
\end{equation}

\paragraph{Practical use.}
$H(X)$ is independent of $q$ and is usually unknown, so BA is most commonly used to \emph{optimize} representations (maximize the bound) rather than to obtain a calibrated numerical estimate of MI.
Maximizing $\E{\log q(x\mid y)}$ corresponds to training a conditional model $q_\theta(x\mid y)$ by maximum likelihood.

\subsubsection*{2) Donsker--Varadhan (DV) bound (MINE)}

The Donsker--Varadhan representation of KL divergence states that for distributions $P,Q$,
\begin{equation}
    \KL{P}{Q} = \sup_{T}\; \Esub{P}{T} - \log \Esub{Q}{e^{T}},
\end{equation}
where the supremum is over measurable functions $T$ for which the expectations exist.
Applying it with $P=p_{X,Y}$ and $Q=p_Xp_Y$ gives
\begin{equation}
\label{eq:dv-bound}
    I(X;Y) \ge I_{\mathrm{DV}}(T)
    \coloneqq \Esub{p(x,y)}{T(x,y)} - \log \Esub{p(x)p(y)}{e^{T(x,y)}}.
\end{equation}
The optimal $T^*$ is $T^*(x,y)=\log r(x,y)+c$ (any constant $c$).

\paragraph{Remarks.}
The DV bound is tight in principle, but in practice it can suffer from high variance because of the $\log\E{e^{T}}$ term.
The estimator known as \emph{MINE} (Mutual Information Neural Estimation) is based on optimizing \eqref{eq:dv-bound} with a neural critic $T_\theta$.

\subsubsection*{3) Nguyen--Wainwright--Jordan (NWJ) bound}

A more numerically stable alternative replaces the log-partition term using the inequality
\(\log u \le u-1\) for $u>0$.
Applying this to \eqref{eq:dv-bound} yields
\begin{equation}
\label{eq:nwj-bound}
    I(X;Y)
    \ge I_{\mathrm{NWJ}}(T)
    \coloneqq \Esub{p(x,y)}{T(x,y)} - \Esub{p(x)p(y)}{e^{T(x,y)-1}}.
\end{equation}
This bound is also tight for the same optimal critic $T^*(x,y)=\log r(x,y)+1$.

\paragraph{Implementation note.}
Both DV and NWJ require estimating expectations under $p_{X,Y}$ and $p_Xp_Y$.
Given a minibatch $\{(x_i,y_i)\}_{i=1}^N$ from the joint, one typically approximates the product-of-marginals expectation by pairing $x_i$ with $y_j$ for $j\neq i$.

\subsection{InfoNCE and Contrastive Learning}
\label{sec:infonce}

InfoNCE is a popular MI lower bound that looks like a multi-class classification objective.
It is the basis of many contrastive methods (SimCLR, MoCo, CLIP, \dots).

\subsubsection*{InfoNCE bound}

Let $f(x,y)$ be a critic score.
Given one positive pair $(x,y)\sim p_{X,Y}$ and $K-1$ negatives $\{y^-_k\}_{k=1}^{K-1}\sim p_Y$ i.i.d., define
\begin{equation}
\label{eq:infonce-loss}
    \mathcal{L}_{\mathrm{InfoNCE}}(f)
    \coloneqq -\E{\log \frac{e^{f(x,y)}}{e^{f(x,y)}+\sum_{k=1}^{K-1} e^{f(x,y^-_k)}}}.
\end{equation}
Then
\begin{equation}
\label{eq:infonce-bound}
    I(X;Y) \ge \log K - \mathcal{L}_{\mathrm{InfoNCE}}(f).
\end{equation}

\paragraph{Two key consequences.}
\begin{itemize}[leftmargin=2em]
    \item \textbf{Upper limit:} the bound cannot exceed $\log K$.
    Therefore, estimating large MI requires many negatives.

    \item \textbf{Classification viewpoint:}
    InfoNCE trains a classifier to identify the matching $y$ among $K$ candidates.
    The optimal critic is (up to an additive $y$-dependent constant) the log conditional density:
    \(f^*(x,y)=\log p(y\mid x)+c(x)\), or equivalently the log density ratio $\log r(x,y)$.
\end{itemize}

\subsubsection*{CLIP-style symmetric InfoNCE}

CLIP learns two encoders $g_\mu(x)\in\fR^d$ and $h_\phi(y)\in\fR^d$ and uses a separable critic
\begin{equation}
    f_{\mu,\phi}(x,y)=\frac{1}{\tau}\, g_\mu(x)^\top h_\phi(y),
\end{equation}
typically with normalized embeddings and a temperature $\tau>0$.
With a minibatch $\{(x_i,y_i)\}_{i=1}^N$, the (one-direction) InfoNCE objective is
\begin{equation}
\label{eq:clip-loss}
    \mathcal{L}_{x\to y}
    = -\frac{1}{N}\sum_{i=1}^N \log \frac{\exp(f(x_i,y_i))}{\sum_{j=1}^N \exp(f(x_i,y_j))}.
\end{equation}
A symmetric loss averages $\mathcal{L}_{x\to y}$ and $\mathcal{L}_{y\to x}$ (swap roles of $x$ and $y$).
This is exactly a pair of cross-entropy losses over the $N\times N$ similarity matrix.

\subsection{What These Bounds Do and Do Not Give You}
\label{sec:mi-est-practice}

\paragraph{Optimization vs. estimation.}
Many ``MI estimators'' are primarily useful as \emph{training objectives}:
maximizing a lower bound encourages dependence.
However, the numerical value of the bound can be a poor approximation of the true MI in high dimensions.
InfoNCE in particular saturates at $\log K$.

\paragraph{Bias/variance tradeoff.}
DV is tight but can be unstable (variance from $\exp(T)$).
NWJ is more stable but can be looser for a fixed critic class.
InfoNCE is very stable and scalable, but biased downward due to the finite number of negatives.

\paragraph{Takeaway.}
Choose the bound based on the goal:
\begin{itemize}[leftmargin=2em]
    \item If you want a \emph{scalable contrastive objective}, use InfoNCE.
    \item If you want a \emph{tighter variational MI lower bound} and can afford careful optimization, consider DV/NWJ-style objectives.
    \item If you already have a good conditional model $q(x\mid y)$, BA can be convenient.
\end{itemize}
